\documentclass[10pt,a4paper]{amsart}
\usepackage[margin=19mm]{geometry}
\usepackage{amsmath,amssymb,mathtools}
\usepackage[scr=rsfs,cal=euler]{mathalfa}
\usepackage{xfrac}
\usepackage[unicode,hidelinks,bookmarks=false]{hyperref}
\newcommand{\ie}{i.e.\ }
\newcommand{\cf}{cf.\ }
\newcommand{\resp}{respectively\ }
\newcommand{\Subsection}{\subsection}
\providecommand{\nobreakdash}{\nobreak\hbox{-}}
\setlength{\emergencystretch}{2em}
\multlinegap0pt
\usepackage{amssymb}
\usepackage{tcolorbox}
\usepackage[shortlabels]{enumitem}
\usepackage{float}
\usepackage{xcolor}
\newtheorem{theorem}{Theorem}
\def\P{{\mathcal P}}
\newcommand{\R}{\mathbb{R}}
\def\U{{\mathcal U}}
\def\eq#1{(\ref{#1})}
\title{On spiral steady flows for the Couette-Taylor problem}
\author{Edoardo Bocchi, Filippo Gazzola, Antonio Hidalgo-Torn\'e}
\date{Source: arXiv:2603.05401v1 (CC BY 4.0)}
\begin{document}
\maketitle

\noindent\emph{Source and licence.} On spiral steady flows for the Couette-Taylor problem. Version arXiv:2603.05401v1 (CC BY 4.0), distributed by arXiv under the Creative Commons Attribution 4.0 International licence, http://creativecommons.org/licenses/by/4.0/, as recorded in the arXiv licence metadata for this exact version and shown on the abstract page https://arxiv.org/abs/2603.05401v1. Full licence notice: \texttt{licenses/arxiv-2603.05401-CC-BY-4.0.txt}.\par\smallskip

\noindent\textit{Editorial scope.} Counted unit: the article's theorem theo-tcg (source0.tex lines 367--374). The article's complete original proof of it is delivered with it (source0.tex lines 375--478, one \texttt{proof} environment of 104 lines). No mathematics is written, repaired or re-derived: the statement and the proof are the article's own lines, in the article's own notation, and no retained line is substituted or re-flowed.  Retained uncounted as the source-local context the proof needs: lines 250--264; lines 267--272; lines 279--283; lines 342--344; lines 360--363.  Nothing was omitted from any retained span, so the package carries no omission record.  No retained line contains a citation, so the wrapper carries no bibliography and no bibliography entry.  No retained line contains a figure environment, an image-inclusion command or drawing code, so no image is delivered and the package declares no asset dependency.  Editorial formatting only, in addition: an AMS \texttt{amsart} preamble that restates the article's own declarations and macros used by the retained lines, namely the environments \texttt{theorem} on separate counters, together with the standard packages those lines need.  The retained environments print in this document as Theorem 1 in order of appearance, whereas the article prints the same items with the numbering of its own full document.  The complete native source of the article as distributed in the arXiv e-print for this exact version is bundled unchanged as \texttt{sources/195798/source0.tex}.
	\begin{equation}\label{SNSEcyl}
		\left\{\begin{array}{l}
			\frac{1}{\rho}\frac{\partial u_\rho}{\partial\rho}+\frac{\partial^2 u_\rho}{\partial\rho^2}+\frac{1}{\rho^2}\frac{\partial^2 u_\rho}{\partial\theta^2}+\frac{\partial^2 u_\rho}{\partial z^2}
			-u_\rho\frac{\partial u_\rho}{\partial\rho}-\frac{u_\theta}{\rho}\frac{\partial u_\rho}{\partial\theta}-u_z\frac{\partial u_\rho}{\partial z}
			+\frac{u_\theta^2}{\rho}-\frac{u_\rho}{\rho^2}-\frac{2}{\rho^2}\frac{\partial u_\theta}{\partial\theta}=\frac{\partial p}{\partial\rho}\\[5pt]
			\frac{1}{\rho}\frac{\partial u_\theta}{\partial\rho}+\frac{\partial^2 u_\theta}{\partial\rho^2}
			+\frac{1}{\rho^2}\frac{\partial^2 u_\theta}{\partial\theta^2}+\frac{\partial^2 u_\theta}{\partial z^2}-u_\rho\frac{\partial u_\theta}{\partial\rho}-\frac{u_\theta}{\rho}\frac{\partial u_\theta}{\partial\theta}
			-u_z\frac{\partial u_\theta}{\partial z}-\frac{u_\rho u_\theta}{\rho}-\frac{u_\theta}{\rho^2}+\frac{2}{\rho^2}\frac{\partial u_\rho}{\partial\theta}
			=\frac{1}{\rho}\frac{\partial p}{\partial\theta}\\[5pt]
			\frac{1}{\rho}\frac{\partial u_z}{\partial\rho}+\frac{\partial^2 u_z}{\partial\rho^2}+\frac{1}{\rho^2}\frac{\partial^2 u_z}{\partial\theta^2}+\frac{\partial^2 u_z}{\partial z^2}
			-u_\rho\frac{\partial u_z}{\partial\rho}-\frac{u_\theta}{\rho}\frac{\partial u_z}{\partial\theta}-u_z\frac{\partial u_z}{\partial z}=\frac{\partial p}{\partial z}\\[5pt]
			\frac{\partial u_\rho}{\partial\rho}+\frac{u_\rho}{\rho}+\frac{1}{\rho}\frac{\partial u_\theta}{\partial\theta}
			+\frac{\partial u_z}{\partial z}=0
		\end{array}\right.
	\end{equation}

	\begin{equation}\label{nsstokes0}
		\begin{aligned}
			& u=\alpha \, e_\theta\quad \text{on}\quad \Gamma_1\, ,\qquad & u =0\quad \text{on}\quad \Gamma_2\, ,\\
			& u=0\quad \text{on}\quad \Gamma_1\, ,\qquad & u =\alpha \, e_\theta\quad \text{on}\quad \Gamma_2\, .
		\end{aligned}
	\end{equation}

	\begin{equation} \label{tc1}
		{\U^C}(\rho) =\frac{R_1}{R_2^2 -R_1^2} \left(\frac{R_2^2}{\rho}-\rho\right), \
		\P^C(\rho) =\frac{1}{2}\left[\frac{R_1}{R_2^2-R_1^2}\right]^{2} \left( \rho^{2} - \frac{R_2^4}{\rho^{2}}
		- 4R_2^2 \log\rho \right)\, 
	\end{equation}

	\begin{equation}\label{assumption}
		u_\rho=u_\rho(\rho,\theta,z)\, ,\quad u_\theta=u_\theta(\rho,z)\, ,\quad u_z=u_z(\rho,\theta)\, ,\quad p=p(\rho,z).
	\end{equation}

    \begin{equation}\label{Ualphabeta}
        \U^P(\rho):=
        \frac{1}{4} \left(\rho^2-R_1^2 - \frac{R_2^2-R_1^2}{\log (R_2/R_1)}\log(\rho/R_1) \right)\qquad \text{for} \quad \rho\in [R_1,R_2].
    \end{equation}

	\begin{theorem}\label{theo-tcg}
		Let $({\U^C},{\P^C})$ be as in \eqref{tc1}, let $\U^P$ be as in \eqref{Ualphabeta}. Then, the spiral Poiseuille flows 
		\begin{equation} \label{tcg}\begin{aligned}
				&\U_{\alpha,\beta}(\rho) := \alpha\U^C(\rho)e_\theta+ \beta\U^P(\rho)e_z,\\[5pt]
				& \P_{\alpha,\beta}(\rho, z) := \alpha^2\P^C(\rho)+ \beta z,\end{aligned}  \qquad \text{for} \quad (\rho,z) \in[R_1,R_2]\times \mathbb{R}\,
		\end{equation}
		are the only partially-invariant solutions  to \eqref{SNSEcyl}-\eqref{nsstokes0}$_1$ as $ \beta$ varies in $\mathbb{R}$.
		\end{theorem}

	\begin{proof} Let $u$ be a partially-invariant solution to \eqref{SNSEcyl}. We divide the proof in two steps. \\
		\underline{Step 1:} we show that the following alternative holds:
		\begin{equation}\label{either}
			\mbox{either }u_\theta\mbox{ does not depend on }z\ (u_\theta=u_\theta(\rho))
			\mbox{ or }u_z\mbox{ does not depend on }\theta\ (u_z=u_z(\rho))\, .
		\end{equation}By \eqref{assumption}, the incompressibility condition becomes
		\begin{equation}\label{incom-cylin}
			\frac{\partial}{\partial\rho}\Big(\rho u_\rho(\rho,\theta,z)\Big)=0 \ \Longrightarrow \ u_{\rho}(\rho,\theta,z)
			=\frac{A(\theta,z)}{\rho} \quad  \forall (\rho,\theta,z)\in(R_1,R_2)\times (0, 2\pi)\times \mathbb{R}\, ,
		\end{equation}
		for some smooth $A:\R^2\to\mathbb{R}$, $2\pi$-periodic with respect to $\theta$. The boundary conditions \eq{nsstokes0}$_1$ imply $A(\theta,z)\equiv0$ and $u_{\rho}(\rho, \theta, z)\equiv0$.
		Therefore, the first three equations in \eqref{SNSEcyl} in the class of partially-invariant solutions reduce to
		\begin{eqnarray}
			\qquad \frac{\partial p}{\partial \rho}(\rho, z) &=& \frac{u^{2}_{\theta}(\rho,z)}{\rho}\, ,\label{uno}\\
			0 \quad &=& \rho\frac{\partial^2 u_{\theta}}{\partial\rho^2}(\rho,z) +
			\frac{\partial u_{\theta}}{\partial\rho}(\rho,z)+\rho\frac{\partial^2 u_{\theta}}{\partial z^2}(\rho,z)
			-\rho u_z(\rho,\theta)\, \frac{\partial u_{\theta}}{\partial z}(\rho,z)- \frac{u_\theta(\rho,z)}{\rho}\, ,\label{dos}\\
			\frac{\partial p}{\partial z}(\rho, \theta) &=& \frac{\partial^2 u_z}{\partial\rho^2}(\rho,\theta) +\frac{1}{\rho}\frac{\partial u_z}{\partial\rho}(\rho,\theta)+\frac{1}{\rho^2}\frac{\partial^2 u_z}{\partial\theta^2}(\rho,\theta)-
			\frac{u_\theta(\rho,z)}{\rho}\, \frac{\partial u_z}{\partial\theta}(\rho,\theta)\, .\label{tres}
		\end{eqnarray}
		From \eqref{dos} we obtain that
		\begin{equation}\label{sinB}
			\rho u_z(\rho,\theta)\, \frac{\partial u_{\theta}}{\partial z}(\rho,z)=\rho\frac{\partial^2 u_{\theta}}{\partial\rho^2}(\rho,z) +
			\frac{\partial u_{\theta}}{\partial\rho}(\rho,z)+\rho\frac{\partial^2u_{\theta}}{\partial z^2}(\rho,z)-\frac{u_\theta(\rho,z)}{\rho}
		\end{equation}
		for all $(\rho, \theta, z)\in (R_1,R_2)\times (0, 2\pi)\times \mathbb{R}$. Since the term on the right-hand side of \eqref{sinB} does not depend on $\theta$ neither does the term on the left, which implies that
		$$
		\text{either}\quad \frac{\partial u_{\theta}}{\partial z}(\rho,z)=0,\quad \text{or}\quad u_z(\rho,\theta)=u_z(\rho),
		$$
		which is equivalent to \eqref{either}.\\
		\underline{Step 2:} we prove that, necessarily, $(u,p)=(\U_{\alpha,\beta}, \P_{\alpha,\beta})$ given by \eqref{tcg}
		for some $\beta\in\R$. To this end, we distinguish two cases.\par
		(a) If the first alternative in \eqref{either} holds, then \eqref{sinB} becomes
		$$
		\rho^2\frac{d^2 u_{\theta}}{d\rho^2}(\rho)+\rho\frac{du_{\theta}}{d\rho}(\rho)-u_\theta(\rho)=0\quad\text{in}\quad (R_1,R_2)
		$$
		which, complemented with the boundary conditions \eq{nsstokes0}$_1$, gives
		\begin{equation}\label{utheta}
			u_\theta(\rho)=\frac{\alpha R_1}{R_2^2 -R_1^2} \left(\frac{R_2^2}{\rho}-\rho\right) .
		\end{equation}
		By replacing \eqref{utheta} into \eqref{uno} and integrating with respect to $\rho$, we find a separated structure of the pressure, namely,
		$
		p(\rho, z)=D(\rho)+E(z)
		$
		for some smooth functions $D:[R_1,R_2]\to\R$ and $E:\R\to\R$. Then, \eqref{tres} reads
		\begin{equation}\label{eq-uz}
			\frac{\partial^2 u_z}{\partial\rho^2}(\rho,\theta)+\frac{1}{\rho}\frac{\partial u_z}{\partial\rho}(\rho,\theta)+\frac{1}{\rho^2}\frac{\partial^2 u_z}{\partial\theta^2}(\rho,\theta)-\frac{u_\theta(\rho)}{\rho}\, \frac{\partial u_z}{\partial\theta}(\rho,\theta)=E'(z)\, .
		\end{equation}
		Since the left hand side of \eqref{eq-uz} does not depend on $z$, neither does the right hand side and $E'(z)\equiv\beta$ for some $\beta\in\R$.
		Thus,
		\begin{equation}\label{sepa-pres}
			p(\rho,z)=D(\rho)+\beta z
		\end{equation}and,
		recalling \eqref{utheta}, \eqref{eq-uz} reduces to
		\begin{equation}\label{F}
			\frac{\partial^2 u_z}{\partial\rho^2}(\rho,\theta)+\frac{1}{\rho}\frac{\partial u_z}{\partial\rho}(\rho,\theta)+
			\frac{1}{\rho^2}\frac{\partial^2 u_z}{\partial\theta^2}(\rho,\theta)+
			\frac{\alpha R_1 }{R_2^2 -R_1^2} \left(1-\frac{R_2^2}{\rho^2}\right)\, \frac{\partial u_z}{\partial\theta}(\rho,\theta)=\beta\in\R\, .
		\end{equation}
		For $2\pi$-periodic functions $\theta\mapsto v(\rho, \theta)$, we restrict ourselves to the class
		\begin{equation}\label{V}
			V=\left\{v\in H^1((R_1,R_2)\times [0,2\pi));\ v(R_1,\theta)=v(R_2, \theta)=0  \text{ for $\theta\in[0,2\pi)$}\right\},
		\end{equation}which is a closed subspace of $H^1((R_1,R_2)\times (0,2\pi))$. 	We claim that \eqref{F} admits a unique solution in $V$.
		Indeed, after multiplying \eqref{F} by $-\rho$, its variational formulation reads
		\begin{equation}\label{var-form}
			B(u_z, v)= F(v) \ \text{for any} \ v\in V,
		\end{equation}
		where $F(v)=-\beta\int_{(R_1,R_2)\times (0,2\pi)} \rho vd\rho d\theta$ is a linear functional on $V$ and \begin{equation}\label{B}B(u_z, v)=\int_{(R_1,R_2)\times (0,2\pi)} \left(\rho \frac{\partial u_z}{\partial \rho}\frac{\partial v}{\partial \rho} + \frac{1}{\rho}\frac{\partial u_z}{\partial \theta} \frac{\partial v}{\partial \theta} + w_\alpha \frac{\partial u_z}{\partial \theta}v\right) d\rho d\theta\end{equation} is a bilinear form on $V$ with $w_\alpha(\rho)= \frac{\alpha R_1}{R_2^2-R_1^2}\left( \frac{R_2^2}{\rho} -\rho\right)$. Thanks to the Hölder inequality, we have that $|F(v)|\leq C_1\|v\|_V$ and $|B(u_z, v)|\leq C_2 \|u_z\|_{V}\|v\|_{V}$ for some constants $C_1=C_1(R_1,R_2,\beta)$ and $C_2=C_2(R_1,R_2, \alpha)$. Furthermore, using the periodicity with respect to $\theta$ in \eqref{V}, we obtain that
		\begin{equation*}\begin{aligned}
				B(v, v)&= \int_{(R_1,R_2)\times (0,2\pi)} \left( \rho\left|\frac{\partial v}{\partial \rho}\right|^2 + \frac{1}{\rho}\left|\frac{\partial v}{\partial \theta}\right|^2 +\frac{1 }{2}\frac{\partial (w_\alpha v^2)}{\partial \theta} \right) d\rho d\theta\\
				&= \int_{(R_1,R_2)\times (0,2\pi)} \left( \rho \left|\frac{\partial v}{\partial \rho}\right|^2 + \frac{1}{\rho}\left|\frac{\partial v}{\partial \theta}\right|^2  \right) d\rho d\theta\geq \min (R_1, 1/R_2)\|v\|^2_V\qquad\forall v\in V,
			\end{aligned}
		\end{equation*} that is, $B$ is coercive (or $V$-elliptic). Then, for any $(\alpha, \beta)\in \mathbb{R}^2$, the Lax-Milgram Theorem yields the existence and uniqueness of the solution to \eqref{F} in $V$. We now seek the solution $u_z$ as a function depending only on $\rho$, hence it satisfies
		\begin{equation*}
			\begin{cases}
				\dfrac{d^2 {u_z}}{d\rho^2}(\rho) +\dfrac{1}{\rho}\dfrac{d{u_z}}{d\rho}(\rho) =\beta \quad \text{in} \quad (R_1,R_2),\\[5pt]
				{u}_{z}(R_1)=  {u}_{z}(R_2)=0,
			\end{cases}
		\end{equation*}which can be solved to obtain the explicit expression
		\begin{equation}\label{uz-explicit}
			{u}_{z}(\rho)=  \frac{\beta}{4} \left(\rho^2-R_1^2 - \frac{R_2^2-R_1^2}{\log (R_2/R_1)}   \log (\rho/R_1)\right).
		\end{equation}
		Thus, by uniqueness, we infer that \eqref{uz-explicit} is the unique solution to \eqref{F}, and we remark that it does not depend on the parameter $\alpha$ (rotational velocity of the inner cylinder $\Gamma_1$).	
		Furthermore, expliciting the function $D(\rho)$ in \eqref{sepa-pres} using \eqref{utheta} yields
		\begin{equation*}
			p(\rho,z)=   \frac{1}{2}\left(\frac{\alpha R_1}{R_2^2-R_1^2} \right)^{2} \left( \rho^{2} - \frac{R_2^4}{\rho^{2}}
			- 4R_2^2 \log\rho \right) + \beta z,
		\end{equation*}
		where we have set the additive constant equal to zero.
		
		(b) If the second alternative in \eqref{either} holds, on the one hand \eqref{tres} reduces to
		$$
		\frac{\partial p}{\partial z}(\rho) = \frac{d^2 u_z}{d\rho^2}(\rho)+\frac{1}{\rho}\frac{d u_z}{d\rho}(\rho)\,,
		$$
		which implies that $\tfrac{\partial^2p}{\partial z\partial\rho}$ merely depends on $\rho$. On the other hand, by differentiating \eqref{uno} with respect to $z$, we obtain
		$$
		\frac{\partial u^2_{\theta}}{\partial z}(\rho,z)=\rho\frac{\partial^2p}{\partial z\partial\rho}(\rho)=: C(\rho)
		$$and integrating with respect to $z$ yields
		\begin{equation}
			| u_\theta (\rho,z)|= \sqrt{C(\rho)z + K(\rho)}
		\end{equation}
		with some smooth function $K: [R_1,R_2] \rightarrow \mathbb{R}$. Since $u_\theta$ is smooth and defined for all $z\in\mathbb{R}$, it must
		necessarily be $C(\rho)\equiv 0$ and we find that $u_\theta=u_\theta(\rho)$; this is the first alternative in \eqref{either}
		that has already been discussed in case (a). The theorem is so proved in both cases.\end{proof}

\end{document}
